\section{Introduction}
\label{sec:introduction}
A coordinate sweep visits each coordinate of a binary position once,
independently swapping or retaining the two positions in every coordinate
pair. A single card is uniform after one sweep. The full deck is a more
difficult object: cards share switches, so their joint permutation law can
retain information that every one-card test misses. We bound that
information in every irreducible representation of the symmetric group.
The bounds imply that an absolute number of sweeps suffices for the full
permutation density to converge to uniform in squared $L^2$ distance as
the deck size tends to infinity.

Let $n=2^d$, $d\ge1$, and identify positions with $\{0,1\}^d$.
A sweep first independently swaps or retains every pair in coordinate
direction 1, and then does the same in directions $2,\ldots,d$.
Write $B_n$ for its probability measure and its average action in a
representation. One sweep is $d$ physical Thorp shuffles, after removing
the deterministic rotations of coordinates; Section~\ref{sec:prelim}
gives the exact convention. For $\lambda\vdash n$, let $V_\lambda$ be
the complex Specht module and $D_\lambda=\dim V_\lambda$. Every trace
in this paper is unnormalized.

\subsection*{Prescribed signed occurrences}
The main estimate describes a representation using two kinds of tensor
colors and a remainder. Permuting tensor factors carrying odd colors
introduces a sign; even
colors describe long rows of a Young diagram, and odd colors describe
long columns. The remainder is represented by distinctly marked positions.
Its size enters the bound with an entropy saving that becomes important
when those marks are sparse.

Precisely, let $u+v+l=n$ and let $\alpha\vdash u$, $\beta\vdash v$,
$\gamma\vdash l$. Suppose
\begin{equation}\label{ten:spin:eq:1}
 V_\alpha\otimes V_{\beta^{\mathrm t}}\otimes V_\gamma
 \subseteq\operatorname{Res}^{S_n}_{S_u\times S_v\times S_l}V_\lambda,
\end{equation}
where $\subseteq$ means occurrence, with no multiplicity-one assumption,
and $\beta^{\mathrm t}$ is the transposed partition. Empty factors are
allowed and have dimension one. Define the unnormalized entropy of the
concatenated part lists by
\begin{equation}\label{ten:spin:eq:2}
 F(\alpha,\beta)=\sum_{i:\alpha_i>0}\alpha_i\log\frac{u+v}{\alpha_i}
 +\sum_{j:\beta_j>0}\beta_j\log\frac{u+v}{\beta_j}.
\end{equation}
It is zero if $u+v=0$. All logarithms are natural.

\begin{theorem}[A moment bound for every signed occurrence]
\label{ten:spin:main}\label{s:spin-moment}
For each sufficiently small fixed $\bar\eta>0$, choose a sufficiently
small $\bar\kappa>0$ with $\bar\kappa<\bar\eta/256$. There is an
integer $r\ge1$, depending only on these fixed constants, such that for
every $d\ge1$, every $\lambda\vdash n=2^d$, and every
occurrence~\eqref{ten:spin:eq:1},
\begin{equation}\label{ten:spin:eq:3}
 \log\!\left(D_\lambda\operatorname{Tr}_{V_\lambda}(B_n^*B_n)^r\right)
 \le\eta_d F(\alpha,\beta)+g_n(l),
\end{equation}
where
\[
 \eta_d=\bar\eta\left(1-\frac1{2\sqrt d}\right),\qquad
 \kappa_d=\bar\kappa\left(1-\frac1{2\sqrt d}\right),\qquad
 g_n(l)=\max\{0,\kappa_d l\log n-l\log(n/l)\}.
\]
Set $g_n(0)=0$ and interpret the logarithm of a zero trace as $-\infty$.
\end{theorem}

The occurrence is prescribed, rather than selected by minimizing the
right-hand side. This is useful in the proof itself: a local restriction
supplies particular signed factors and a particular remainder, and the
same theorem applies to those data. The saving $-l\log(n/l)$ is also
essential. The two child bounds retain enough of it to pay for placing
the marks when the children are assembled into the parent rectangle.

\subsection*{The full permutation law}
For a fixed number $L$ of sweeps, let
$f_{d,L}(\sigma)=n!\,B_n^{*L}(\sigma)$ be the density of their
permutation law with respect to the uniform measure $U_{S_n}$.
Let $t_{\mathrm{mix}}(d)$ be the least number of physical shuffles
whose law has total variation distance at most $1/4$ from uniform.

\begin{corollary}[Fixed-sweep convergence of the full deck]
\label{cor:spin-full-deck}
There is an absolute positive integer $L$ such that
\[
 \mathbb E_{U_{S_n}}|f_{d,L}-1|^2\longrightarrow0
 \qquad(d\longrightarrow\infty).
\]
The convergence is uniform over starting orders. Moreover
$t_{\mathrm{mix}}(d)=\Theta(d)$.
\end{corollary}

To obtain the upper bound, choose an occurrence with no remainder and
$F(\alpha,\beta)\le C_0\log D_\lambda$. With $\bar\eta$ sufficiently
small, the theorem then gives a negative power of $D_\lambda$ for the
sweep norm. Summing those powers by finite-group Fourier analysis proves
the corollary. Section~\ref{ten:spin} establishes the required occurrence
and completes this argument. The lower bound comes from the number of
available switch choices: $t$ physical shuffles have support of size at
most $2^{tn/2}$, whereas the uniform law is supported on $n!$ permutations.
Lemma~\ref{lem:support} gives $t_{\mathrm{mix}}(d)\ge2d-O(1)$. The
same fixed-sweep estimate also gives an approximately uniform permutation
sampler whose output coordinates each require only $O(\log n)$ adaptive
queries to shared randomness; see Corollary~\ref{cor:decision-forest}.

\subsection*{The proof mechanism}
There are two sources of contraction. If a representation first appears
on a small number of tracked cards, centering the tuple kernel cancels
every collection of paths with an isolated member. A weighted forest
count bounds the remaining encounters. For the other representations,
split the coordinates into two nearly equal groups. Positions form a
rectangle: the first sub-sweep acts within columns, the second within
rows. The corresponding representation decompositions do not commute,
so the proof must bound the angle between their tensor type spaces.

The angle estimate uses a positive matrix of trace one on each line.
Averaging its tensor powers over the parity-preserving unitary group
dominates the projection to a prescribed signed type. The mean column
density supplies an inverse square root that normalizes all the cells
in a row. The resulting one-cell factors have mean at most one on the
complete rectangle. Removing $l$ cells costs at most $e^l$, while the
global type retains its entropy saving. Section~\ref{sec:static} proves
this comparison, including the local carrier dimensions, before any
moment induction is used.

Distinct marks supply the other local input. As Figure~\ref{fig:marked-grid}
shows, a row move followed by a column move can return a named point
to its initial cell only if both moves fix that point. The trace therefore restricts group-algebra
coefficients to a pointwise stabilizer. This restriction is positive,
and its exact regular-trace factor is the reciprocal of a falling
factorial. Summing mark placements retains the entropy term needed by
the parent bound.

\begin{figure}[t]
\centering
\begin{tikzpicture}[x=.64cm,y=.64cm]
\fill[blue!14] (0,2) rectangle (7,3);
\fill[orange!20] (2,0) rectangle (3,4);
\fill[green!20!white] (2,2) rectangle (3,3);
\draw[step=1,gray!65] (0,0) grid (7,4);
\draw[blue!70!black,thick] (0,2) rectangle (7,3);
\draw[orange!75!black,thick] (2,0) rectangle (3,4);
\node at (2.5,2.5) {$a$};
\node[left] at (0,2.5) {$i$};
\node[above] at (2.5,4) {$j$};
\node[blue!70!black,anchor=west,align=left] at (7.4,3)
  {after a row move:\\same row $i$};
\node[orange!75!black,anchor=west,align=left] at (7.4,1)
  {before a column return:\\same column $j$};
\draw[blue!70!black,-{Stealth}] (7.35,2.5)--(6.5,2.5);
\draw[orange!75!black,-{Stealth}] (7.35,.5)--(2.5,.5);
\node[align=center] at (3.5,-.65)
  {the intermediate cell lies in their intersection: $(i,j)$};
\end{tikzpicture}

\caption{One distinct mark initially occupies cell $(i,j)$ in this
schematic rectangle. After a row move its intermediate location lies
in the highlighted row; a column move can return it to $(i,j)$ only
from the highlighted column. Their intersection is the initial cell.
The same argument for every distinct label forces both moves to fix
all marks. Unmarked cells carry signed tensor factors.}
\label{fig:marked-grid}
\end{figure}

\subsection*{History and mathematical background}
The Thorp shuffle arose in the study of shuffling procedures associated
with Faro~\cite{Thorp1973}. Its simple local rule long resisted a
full-deck mixing analysis. For $n=2^d$, Morris's first polynomial
bound was $O(d^{44})$, using the chameleon process and evolving
sets~\cite{Morris2008}; Montenegro and Tetali improved this to
$O(d^{29})$ through spectral-profile and evolving-set
estimates~\cite[Section~6.4]{MontenegroTetali2006}.
Using entropy contraction, Morris later proved an $O((\log n)^4)$ bound for all even deck
sizes~\cite{Morris2009}, followed by $O(d^3)$ for power-of-two
decks~\cite{Morris2013}. All these bounds count physical shuffles,
rather than sweeps of $d$ coordinate layers. Corollary~\ref{cor:spin-full-deck}
gives $O(d)$ physical shuffles and convergence of the full permutation
density in squared $L^2$, with the matching order supplied by the
support lower bound.

Partial permutation laws are also central to the cryptographic study
of Thorp shuffling~\cite{mrs,MorrisRogawayStegers2018}. Such results
can track many cards: Czumaj and V\"ocking proved an
$O((\log n)^2)$ bound for the joint positions of any fixed fraction
$cn$ of the cards, where $0<c<1$, using non-Markovian
coupling~\cite{CzumajVocking2014}. The distinction here is the passage
to the joint law of all $n$ cards. The signed representation estimates
retain that complete permutation information instead of selecting a
prescribed family of card marginals.

We use the classical branching and Littlewood--Richardson rules, the
hook-length formula, and ordinary Schur--Weyl
duality~\cite{James1978,Sagan2001,Stanley1999,etingof,VershikOkounkov2005}.
Signed tensor and hook combinatorics have their representation-theoretic
ancestry in Berele and Regev; hook-Schur formulas and their positive
tableau descriptions are also developed by Orellana and Zabrocki and
by Mason and Niese~\cite{BereleRegev1987,orellana-zabrocki2000,mason-niese2016}.
Our decomposition uses ordinary Schur--Weyl duality on the two parity
classes, a sign twist, and induction. The density comparisons and the
budget estimates are proved here.

The passage from Fourier estimates to total variation is the
finite-group method of Diaconis and
Shahshahani~\cite{DiaconisShahshahani1981}. Their random-transposition
law is central. A coordinate sweep is generally nonnormal: its product
with its adjoint need not equal the product in the opposite order.
We use the Araki--Lieb--Thirring inequality for positive
matrices~\cite{araki1990,audenaert2007} to recurse on moments, and
Schatten H\"older to pass to repeated forward sweeps.

\subsection*{The later bounds}
Theorem~\ref{s:spin-moment} and Corollary~\ref{cor:spin-full-deck} are
proved in Section~\ref{ten:spin}. The remainder of the paper studies
what additional information a moment estimate can retain. Some bounds
optimize over signed factors and marks; others keep every prescribed
hook subdiagram, minimize over all subdiagrams, or couple a first-row
defect bound with a row-and-column weight. Their statements differ in
both their optimization domains and their moment exponents.

These later sections also develop distinct mechanisms: lowering
operators give a norm estimate throughout $2\le k\le n/2$;
covariance retains local ranks; a harmonic-tuple coupling gives a
separate sparse estimate; and ordered products of densities provide
another marked-grid proof. They supply alternative routes to the same
mixing order, as well as their stated finer estimates.
Part~\ref{part:refinements} begins with the guide in Section~\ref{s:later-routes}, which
compares these targets and the information each retains.
Section~\ref{s:full-isotypic-section} then provides the common
full-isotypic, positive-mixture, coefficient-integral, and
marked-recurrence tools before the individual arguments begin.
